% Options for packages loaded elsewhere
\PassOptionsToPackage{unicode}{hyperref}
\PassOptionsToPackage{hyphens}{url}
\PassOptionsToPackage{dvipsnames,svgnames,x11names}{xcolor}
\documentclass[
  10pt,
]{article}
\usepackage{xcolor}
\usepackage[margin=2.4cm]{geometry}
\usepackage{amsmath,amssymb}
\setcounter{secnumdepth}{-\maxdimen} % remove section numbering
\usepackage{iftex}
\ifPDFTeX
  \usepackage[T1]{fontenc}
  \usepackage[utf8]{inputenc}
  \usepackage{textcomp} % provide euro and other symbols
\else % if luatex or xetex
  \usepackage{unicode-math} % this also loads fontspec
  \defaultfontfeatures{Scale=MatchLowercase}
  \defaultfontfeatures[\rmfamily]{Ligatures=TeX,Scale=1}
\fi
\usepackage{lmodern}
\ifPDFTeX\else
  % xetex/luatex font selection
  \setmainfont[]{Libertinus Serif}
  \setmonofont[Scale=0.85]{Latin Modern Mono}
  \setmathfont[]{Latin Modern Math}
\fi
% Use upquote if available, for straight quotes in verbatim environments
\IfFileExists{upquote.sty}{\usepackage{upquote}}{}
\IfFileExists{microtype.sty}{% use microtype if available
  \usepackage[]{microtype}
  \UseMicrotypeSet[protrusion]{basicmath} % disable protrusion for tt fonts
}{}
\makeatletter
\@ifundefined{KOMAClassName}{% if non-KOMA class
  \IfFileExists{parskip.sty}{%
    \usepackage{parskip}
  }{% else
    \setlength{\parindent}{0pt}
    \setlength{\parskip}{6pt plus 2pt minus 1pt}}
}{% if KOMA class
  \KOMAoptions{parskip=half}}
\makeatother
\usepackage{longtable,booktabs,array}
\newcounter{none} % for unnumbered tables
\usepackage{calc} % for calculating minipage widths
% Correct order of tables after \paragraph or \subparagraph
\usepackage{etoolbox}
\makeatletter
\patchcmd\longtable{\par}{\if@noskipsec\mbox{}\fi\par}{}{}
\makeatother
% Allow footnotes in longtable head/foot
\IfFileExists{footnotehyper.sty}{\usepackage{footnotehyper}}{\usepackage{footnote}}
\makesavenoteenv{longtable}
\setlength{\emergencystretch}{3em} % prevent overfull lines
\providecommand{\tightlist}{%
  \setlength{\itemsep}{0pt}\setlength{\parskip}{0pt}}

\providecommand{\E}{\mathbb{E}}
\providecommand{\N}{\mathbb{N}}
\providecommand{\Z}{\mathbb{Z}}
\providecommand{\R}{\mathbb{R}}
\providecommand{\eps}{\varepsilon}
\providecommand{\ceil}[1]{\left\lceil #1\right\rceil}
\providecommand{\floor}[1]{\left\lfloor #1\right\rfloor}
\AtBeginDocument{\let\setminus\smallsetminus}
\usepackage[most]{tcolorbox}
\newtcolorbox{warning}{colback=red!5,colframe=red!65!black,boxrule=0.8pt,arc=2pt,left=6pt,right=6pt,top=5pt,bottom=5pt}
\usepackage{etoolbox}
\AtBeginEnvironment{longtable}{\small}
\setlength{\emergencystretch}{3em}
\usepackage{fancyhdr}
\pagestyle{fancy}
\fancyhf{}
\fancyhead[L]{\small\bfseries UNREFEREED GPT-6 ASTRA OUTPUT}
\fancyhead[R]{\small\thepage}
\renewcommand{\headrulewidth}{0.4pt}
\usepackage{bookmark}
\IfFileExists{xurl.sty}{\usepackage{xurl}}{} % add URL line breaks if available
\urlstyle{same}
\hypersetup{
  pdftitle={Conjecture 7},
  pdfauthor={Writeup: gpt-6-astra (single model pass)},
  colorlinks=true,
  linkcolor={blue!50!black},
  filecolor={Maroon},
  citecolor={Blue},
  urlcolor={blue!50!black},
  pdfcreator={LaTeX via pandoc}}

\title{Conjecture 7}
\usepackage{etoolbox}
\makeatletter
\providecommand{\subtitle}[1]{% add subtitle to \maketitle
  \apptocmd{\@title}{\par {\large #1 \par}}{}{}
}
\makeatother
\subtitle{Unrefereed candidate proof by GPT-6 Astra}
\author{Writeup: \texttt{gpt-6-astra} (single model pass)}
\date{Catalog id \texttt{2603.02786\_\_04} --- generated 2026-09-17}

\begin{document}
\maketitle

\begin{warning}

\textbf{UNREFEREED MODEL OUTPUT.} This document was selected solely
because GPT-6 Astra labelled its own result
\texttt{would\_publish:\ true} and returned \texttt{proved} or
\texttt{disproved}. It has not passed the adversarial LLM referee used
for the earlier Sol campaign, has not been checked by a human
mathematician, and has not been checked for novelty. Treat every
mathematical and bibliographic claim below as unverified.

\end{warning}

\section{Summary and provenance}\label{summary-and-provenance}

{\def\LTcaptype{none} % do not increment counter
\begin{longtable}[]{@{}
  >{\raggedright\arraybackslash}p{(\linewidth - 2\tabcolsep) * \real{0.5000}}
  >{\raggedright\arraybackslash}p{(\linewidth - 2\tabcolsep) * \real{0.5000}}@{}}
\toprule\noalign{}
\endhead
\bottomrule\noalign{}
\endlastfoot
Catalog id & \texttt{2603.02786\_\_04} \\
Catalog entry &
\href{https://graph-theory-ai.github.io/graph-conjectures/arxiv/2603.02786__04/}{Conjecture
7} \\
Source corpus & arxiv \\
Campaign leg & selected second attempts (\texttt{attacks\_retry}) \\
Source paper / entry & Packing arithmetic progressions \\
Model verdict & \textbf{proved} (confidence: high) \\
Model's one-line claim & A random-composition fractional tiling, rounded
by edge-colouring and a reservoir, proves M\_k(n)=nk+o(k) for every
fixed n. \\
Model & \texttt{gpt-6-astra}, reasoning effort \texttt{max},
\texttt{mode=pro}, flex service tier \\
Original artifact &
\texttt{attacks\_retry/2603.02786\_\_04/output.md} \\
Independent review & \textbf{None} \\
Model's caveats & Uses the classical Pippenger--Spencer edge-colouring
theorem; no explicit convergence rate is obtained. \\
\end{longtable}
}

\section{Problem statement}\label{problem-statement}

\emph{No extracted statement was stored in the artifact metadata.}

\section{Astra writeup}\label{astra-writeup}

\emph{The text below is the model output verbatim, apart from moving its
machine-readable verdict block into the summary above and shifting
Markdown heading levels for this document. Mechanical TeX defects and
equation tags were normalized where needed for compilation.}

\subsection{Theorem}\label{theorem}

For every fixed positive integer \(n\), \[
M_k(n)=nk+o(k)\qquad(k\to\infty).
\]

Here a translate of the progression with difference \(d\) is written \[
A(d,a)=\{a,a+d,\ldots,a+(n-1)d\}.
\] Thus \(M_k(n)\) is the least \(m\) such that one can choose pairwise
disjoint sets \(A(d,a_d)\subseteq[m]\), one for every \(d\in[k]\).

The missing ingredient in the supplied attempt is an appropriate
non-uniform fractional packing. Such a fractional packing can be
constructed \textbf{exactly}, using a random composition of an integer.
After a slight enlargement and removal of the smallest differences, its
weights satisfy the hypotheses needed for rounding.

The case \(n=1\) is immediate, so below \(n\ge2\) is fixed.

\subsubsection{1. An exact fractional packing from random
compositions}\label{an-exact-fractional-packing-from-random-compositions}

We first construct fractional packings on the optimal number of points.

\paragraph{Lemma 1}\label{lemma-1}

For every positive integer \(K\), there are nonnegative weights \[
u_{d,a},
\qquad
1\le d\le K,\quad 1\le a\le nK-(n-1)d,
\] with the following properties:

\begin{enumerate}
\def\labelenumi{\arabic{enumi}.}
\item
  For every \(d\in[K]\), \[
  \sum_a u_{d,a}=1.
  \qquad\text{(1)}
  \]
\item
  For every \(x\in[nK]\), \[
  \sum_{\substack{d,a\\x\in A(d,a)}}u_{d,a}=1.
  \qquad\text{(2)}
  \]
\item
  For every \(d,a\), \[
  u_{d,a}\le \frac1d.
  \qquad\text{(3)}
  \]
\item
  If \(k<K\le2k\), then, for every \(x\in[nK]\), \[
  \sum_{\substack{d\le k,\ a\\x\in A(d,a)}}u_{d,a}
  \le 1-\frac{K-k}{2K}.
  \qquad\text{(4)}
  \]
\end{enumerate}

\paragraph{Proof}\label{proof}

Generate a random composition of \(K\) as follows. When the remaining
sum is \(r>0\), choose the next part uniformly from \(\{1,\ldots,r\}\),
subtract it, and continue.

Let \(C_d(K)\) be the number of parts of size \(d\). We claim that \[
\mathbb E C_d(K)=\frac1d
\qquad(1\le d\le K).
\qquad\text{(5)}
\] Indeed, writing \(f_r=\mathbb E C_d(r)\), we have \(f_r=0\) for
\(r<d\), while, for \(r\ge d\), \[
f_r=\frac1r+\frac1r\sum_{s=0}^{r-1}f_s.
\] Induction gives \[
f_r=\frac1r+\frac{r-d}{rd}=\frac1d.
\]

Suppose the generated composition is \[
K=d_1+\cdots+d_\ell.
\] Partition \([nK]\), from left to right, into consecutive blocks of
lengths \[
nd_1,\ldots,nd_\ell.
\] A block of length \(nd\), beginning immediately after position \(b\),
is tiled by the following \(d\) disjoint progressions: \[
A(d,b+t)=\{b+t,b+t+d,\ldots,b+t+(n-1)d\},
\qquad 1\le t\le d.
\] Thus the composition produces a tiling of \([nK]\) by \(n\)-term
arithmetic progressions. It contains exactly \(dC_d(K)\) tiles of
difference \(d\).

Finally, independently with probability \(1/2\), reflect the entire
tiling by \[
x\longmapsto nK+1-x.
\] Reflection preserves every tile's difference.

Define \(u_{d,a}\) to be the probability that \(A(d,a)\) is a tile in
the resulting random tiling. Equation (5) gives \[
\sum_a u_{d,a}
=\mathbb E[dC_d(K)]
=1,
\] proving (1). Every outcome tiles every point exactly once, proving
(2).

The appearance of a particular tile of difference \(d\) requires
\(C_d(K)\ge1\). Hence \[
u_{d,a}\le\mathbb P(C_d(K)\ge1)
\le\mathbb E C_d(K)=\frac1d,
\] which proves (3).

It remains to prove (4). The first part \(U\) of the composition is
uniform on \([K]\). If \(U>k\), the first block has length \[
nU\ge n(k+1)>\frac{nK}{2}.
\] Consequently:

\begin{itemize}
\tightlist
\item
  for a point in the left half of \([nK]\), the event ``no reflection
  and \(U>k\)'' guarantees that it lies in a tile of difference greater
  than \(k\);
\item
  for a point in the right half, the event ``reflection and \(U>k\)''
  gives the same guarantee.
\end{itemize}

Each event has probability \[
\frac12\frac{K-k}{K}.
\] Thus, at every point, the total weight belonging to differences
greater than \(k\) is at least \((K-k)/(2K)\). Subtracting this from (2)
proves (4). \(\square\)

The cancellation behind the construction is worth emphasizing: a part of
size \(d\) occurs an expected \(1/d\) times and contributes exactly
\(d\) tiles of difference \(d\).

\subsubsection{2. The edge-colouring
input}\label{the-edge-colouring-input}

We use the following established consequence of the Pippenger--Spencer
asymptotic edge-colouring theorem.

\begin{quote}
\textbf{Edge-colouring theorem.} Fix \(r\ge2\). If \(H\) is an
\(r\)-uniform hypergraph with maximum degree \(\Delta\to\infty\) and
maximum pair-codegree \(o(\Delta)\), then \[
\chi'(H)\le(1+o(1))\Delta.
\qquad\text{(6)}
\] Equivalently, its edges can be partitioned into at most
\((1+o(1))\Delta\) matchings.
\end{quote}

For completeness, the standard regular form of the theorem implies
precisely this maximum-degree formulation. Here is a regularization that
avoids any hidden almost-regularity assumption.

Choose a prime \(p>\max\{r,\Delta\}\), and take copies of \(H\) indexed
by \[
(i,b)\in\{0,\ldots,r-1\}\times\mathbb F_p.
\] Write a copied vertex as \((v,i,b)\). For every \(v\in V(H)\), choose
a set \[
S_v\subseteq\mathbb F_p,
\qquad |S_v|=\Delta-\deg_H(v).
\] For each \(s\in S_v\) and \(b\in\mathbb F_p\), add the edge \[
\{(v,i,b+is):0\le i<r\}.
\qquad\text{(7)}
\] Every copied vertex now has degree exactly \(\Delta\). Two vertices
with different original identities acquire no additional common edge.
Two distinct copies of the same original vertex lie together in at most
one added edge: their two coordinates determine \(s\) and \(b\)
uniquely. Thus the new hypergraph is simple, \(\Delta\)-regular, and has
maximum pair-codegree at most \[
\max\{\Delta_2(H),1\}.
\] Apply the regular Pippenger--Spencer theorem and restrict its
edge-colouring to one original copy of \(H\). This gives (6).

\subsubsection{3. Rounding a spread fractional packing with
slack}\label{rounding-a-spread-fractional-packing-with-slack}

The reservoir argument suggested in the supplied attempt is valid. We
give the needed weighted version, including the concentration and
codegree checks.

\paragraph{Lemma 2}\label{lemma-2}

Fix \(n\ge2\) and constants \(\alpha,C>0\), \(0<\delta<1\). Let
\(S\subseteq[k]\), let \(m=O(k)\), and put \[
L_d=m-(n-1)d.
\] Suppose \(L_d\ge\alpha k\) for every \(d\in S\), and suppose there
are nonnegative weights \(w_{d,a}\) on all admissible translates
\(A(d,a)\subseteq[m]\), \(d\in S\), satisfying \[
\sum_{a=1}^{L_d}w_{d,a}=1
\qquad(d\in S),
\qquad\text{(8)}
\] \[
w_{d,a}\le\frac Ck,
\qquad\text{(9)}
\] and \[
\sum_{\substack{d\in S,\ a\\x\in A(d,a)}}w_{d,a}
\le1-\delta
\qquad(x\in[m]).
\qquad\text{(10)}
\] Then, for all sufficiently large \(k\), all the differences in \(S\)
can be packed in \([m]\).

All asymptotic statements here keep \(n,\alpha,C,\delta\) and the
constant in \(m=O(k)\) fixed.

\paragraph{Proof}\label{proof-1}

We may assume \(S\ne\varnothing\).

Choose a constant \[
0<\rho<\delta,
\qquad \sigma=1-\rho.
\] Independently put each point of \([m]\) into a reservoir \(R\) with
probability \(\rho\), and let \(W=[m]\setminus R\).

Set \[
D_0=\frac{k}{2C}.
\] Independently for every admissible labelled progression, retain it
with probability \[
p_{d,a}=D_0w_{d,a}\le\frac12.
\]

Form an \((n+1)\)-uniform hypergraph \(H\) with:

\begin{itemize}
\tightlist
\item
  one formal label vertex for every \(d\in S\);
\item
  the point vertices of \(W\);
\end{itemize}

and with an edge \[
\{d\}\cup A(d,a)
\] whenever \(A(d,a)\) was retained and is wholly contained in \(W\).

Let \[
D=D_0\sigma^n.
\] For every label \(d\), equation (8) gives \[
\mathbb E\deg_H(d)=D_0\sigma^n=D.
\qquad\text{(11)}
\] For every point \(x\), conditional on \(x\in W\), equation (10) gives
\[
\mathbb E[\deg_H(x)\mid x\in W]
\le D_0\sigma^{n-1}(1-\delta)
=D\frac{1-\delta}{\sigma}.
\qquad\text{(12)}
\] Since \(\rho<\delta\), the last ratio is strictly less than \(1\).

We verify the bounded-codegree and concentration facts used below.

For a label \(d\) and a point \(x\), there are at most \(n\) candidate
edges containing both. For two point vertices \(x<y\), a candidate edge
containing both must satisfy \[
y-x=qd
\] for some \(q\in\{1,\ldots,n-1\}\). For a fixed \(q\), this determines
\(d\), and there are at most \(n-q\) possible starts. Hence,
deterministically, \[
\Delta_2(H)
\le \max\left\{n,\sum_{q=1}^{n-1}(n-q)\right\}
=O_n(1).
\qquad\text{(13)}
\]

For a fixed label degree, changing the reservoir status of one point
changes the degree by at most \(n\). For a fixed point degree, after
conditioning on that point belonging to \(W\), changing another point's
status changes the degree by at most the bound in (13). Changing one
retention coin changes a degree by at most one. Each of these degrees
depends on only \(O(k)\) relevant independent variables.

The bounded-differences inequality, with deviation \(k^{2/3}\),
therefore shows that, simultaneously for all labels and point vertices,
with probability tending to one, \[
\deg_H(d)=(1+o(1))D
\qquad(d\in S),
\qquad\text{(14)}
\] and \[
\deg_H(x)\le(1-\gamma)D
\qquad(x\in W)
\qquad\text{(15)}
\] for some fixed \(\gamma>0\). Indeed, each failure probability is \[
\exp(-\Omega(k^{1/3})),
\] and there are only \(O(k)\) degrees to consider.

We also require a reservoir property. For \(d\in S\), let \[
Y_d=\bigl|\{a:A(d,a)\subseteq R\}\bigr|.
\] Then \[
\mathbb EY_d=\rho^nL_d\ge\rho^n\alpha k.
\] Changing the reservoir status of one point changes \(Y_d\) by at most
\(n\). Another application of bounded differences shows that,
simultaneously for every \(d\in S\), with probability tending to one, \[
Y_d\ge \frac12\rho^n\alpha k=:ck,
\qquad\text{(16)}
\] where \(c>0\) is fixed.

Fix a realization satisfying (14)--(16). By (14)--(15), \[
\Delta(H)=(1+o(1))D,
\] and (13) gives \(\Delta_2(H)=o(D)\). The edge-colouring theorem
partitions \(E(H)\) into at most \[
(1+o(1))D
\] matchings.

Writing \(N=|S|\), every edge contains exactly one label, so \[
|E(H)|=\sum_{d\in S}\deg_H(d)=(1+o(1))ND.
\] Some colour class is therefore a matching containing at least \[
(1-o(1))N
\] edges. It packs all but a set \(T\) of labels, where \[
|T|=o(k).
\]

Pack the remaining labels greedily inside \(R\). For any new difference
\(d\), a previously used point forbids at most \(n\) candidate starts.
Thus one previously selected progression forbids at most \(n^2\)
candidates. Throughout the completion, fewer than \[
n^2|T|=o(k)
\] candidates are forbidden, whereas (16) supplies at least \(ck\)
candidates for every remaining label. The greedy completion succeeds for
sufficiently large \(k\).

The first matching lies in \(W\), while the completion lies in \(R\), so
the two packings are disjoint. \(\square\)

\subsubsection{\texorpdfstring{4. The small differences cost only
\(O_n(s)\)}{4. The small differences cost only O\_n(s)}}\label{the-small-differences-cost-only-o_ns}

We will put an arbitrarily small initial segment of differences in a
separate interval.

For every \(s\ge1\), \[
M_s(n)\le (n^2+n-1)s.
\qquad\text{(17)}
\] To see this, set \[
m=(n^2+n-1)s
\] and greedily place the differences \(1,\ldots,s\). Every difference
\(d\le s\) has at least \[
m-(n-1)d\ge n^2s
\] admissible starts. Each previously used point forbids at most \(n\)
starts, so after at most \(s-1\) progressions have been placed, at most
\[
n^2(s-1)
\] starts are forbidden. There is always an available choice.

\subsubsection{5. Proof of the theorem}\label{proof-of-the-theorem}

Fix a constant \[
0<\eta<\frac12.
\] For sufficiently large \(k\), put \[
K=\lceil(1+\eta)k\rceil,
\qquad
q=\lfloor\eta k\rfloor,
\qquad
m=nK.
\] In particular, \[
k<K\le2k.
\]

Take the weights from Lemma 1 on \([nK]\), but retain only labels \[
S=\{q+1,\ldots,k\}.
\] For \(d\in S\), set \(w_{d,a}=u_{d,a}\).

We check every hypothesis of Lemma 2:

\begin{itemize}
\tightlist
\item
  The label sums remain exactly one, by (1).
\item
  By (3), \[
  w_{d,a}\le\frac1d
  \le\frac1{q+1}
  \le\frac1{\eta k}.
  \qquad\text{(18)}
  \]
\item
  Removing the labels greater than \(k\), and then also those at most
  \(q\), gives by (4) \[
  \sum_{\substack{d\in S,\ a\\x\in A(d,a)}}w_{d,a}
  \le1-\frac{K-k}{2K}
  \le1-\frac{\eta}{4}.
  \qquad\text{(19)}
  \]
\item
  Every retained difference has at least \[
  L_d=nK-(n-1)d
  \ge nK-(n-1)k
  \ge k
  \qquad\text{(20)}
  \] admissible starts.
\item
  Finally, \(m=nK\le2nk=O(k)\).
\end{itemize}

Thus Lemma 2 applies with, for example, \[
\alpha=1,\qquad C=\frac1\eta,\qquad \delta=\frac\eta4.
\] It packs all differences \(q+1,\ldots,k\) in \([nK]\), once \(k\) is
sufficiently large.

By (17), the remaining differences \(1,\ldots,q\) can be packed in a
separate interval of length \[
(n^2+n-1)q.
\] Concatenating the two intervals gives \begin{equation*}
\begin{aligned}
M_k(n)
&\le nK+(n^2+n-1)q\\
&\le nk+(n^2+2n-1)\eta k+n.
\end{aligned}
\qquad\text{(21)}
\end{equation*}

Since \(\eta>0\) was arbitrary, first letting \(k\to\infty\) and then
\(\eta\downarrow0\) yields \[
\limsup_{k\to\infty}\frac{M_k(n)}k\le n.
\] The \(k\) disjoint progressions contain \(nk\) points, so \[
M_k(n)\ge nk.
\] Consequently, \[
\boxed{M_k(n)=nk+o(k)=(1+o(1))nk}
\] for every fixed positive integer \(n\). \(\square\)

\subsubsection{Scope and dependencies}\label{scope-and-dependencies}

The only non-elementary input is the established Pippenger--Spencer
edge-colouring theorem; its applicability, including the maximum-degree
formulation, was checked above. The argument leaves no unresolved
fractional-packing or completion step.

The proof is asymptotic for \textbf{fixed} \(n\). It does not assert
uniformity when \(n\) grows with \(k\), nor does it extract an explicit
rate for the \(o(1)\) term.

\end{document}
